\section{Results and precise comparison}
\label{sec:results}
The principal deliverable is the fixed ordinary polynomial
\begin{equation}\label{eq:principal}
 \boxed{P_{217}(x_0,\ldots,x_{216})
       =\sum_{j=0}^{153}R_j(x_0,\ldots,x_{216})^2=0,
       \qquad(x_0,\ldots,x_{216})\in\Z^{217}.}
\end{equation}
Every residual is printed in Appendix~\ref{sec:residuals}. The literal
arrays and fully collected expansion contain no unspecified parameter,
variable exponent, sequence operator or hidden existential coordinate.
The polynomial has exact degree four, 7396 nonzero monomials and
coefficient height 8880.

\begin{theorem}\label{thm:main}
For the displayed principal polynomial, and for each separately specified
companion on its stated domain $D$, the new ordinary witness arguments
composed with the identified inherited arithmetic and syntactic
calibration give
\[
 \mathrm{ZFC}\vdash
 \bigl[\neg\exists\boldsymbol x\in D^n\;P(\boldsymbol x)=0\bigr]
 \longleftrightarrow\Con(\mathrm{ZFC}).
\]
No consistency assumption is used to prove this equivalence.
\end{theorem}
The assertion is not a proof of either consistency or unsolvability.
The report separates the new unbounded arguments, exact finite algebra,
and the retained logical interfaces. It does not claim a new complete
proof-assistant kernel endpoint or a global minimum.

\subsection{Direct quartics with moderate literal coefficients}
\begin{center}\small
\begin{tabular}{@{}lrrrr@{}}\toprule
Equation over unrestricted integers & Variables & Rows & Terms & Height\\\midrule
Round 20 balanced &238&175&5236&3968\\
Incoming Strict Horizons principal &236&173&4876&7968\\
Incoming Pooled Powers &225&162&7867&128446\\\midrule
New principal &217&154&7396&8880\\
New lower-height companion &217&154&14793&3200\\
New dimension companion &216&153&25602&5632\\\bottomrule
\end{tabular}
\end{center}
All entries have exact degree four. The new principal improves every
listed measure against the strongest new pooled input: eight fewer
variables and rows, 471 fewer terms and lower coefficient height.
Against its actual factorial-free Strict Horizons source it saves
19 variables and rows, but increases support and height. The separate
lower-height and dimension companions express real tradeoffs.
The earlier Round~20 balanced point also retains its sparse, low-height
advantage. There is no single definition-independent meaning of simplest.

The pooled input supplies the simultaneous exact-power module; it is not
the arithmetic ancestor of the new principal. The latter's source chain
runs through Strict Horizons and its factorial-free C222 construction.
The emitted source edits, common-base groups, radical projections and
inverse coordinate maps establish the composition. A coordinate count
obtained merely by adding savings from unrelated packets would not.

\subsection{Height one and compact scalar equations}
At coefficient height one, unit-homogenized equality cycles and shared
coefficient factors reduce support without changing the six selected
variable counts. The 505-variable point goes from 20491 to 20405 terms;
the 503-variable alternative goes from 24193 to 24183. The complete
six-point table and the zero-set proof appear in Section~\ref{sec:signed-network}.

The ten-positive-variable branch supplies three separately fixed
circuits: exact degree/cost $240/82$, $232/85$ and $214/87$.
The degree-214 equation combines a free quotient with inverse-Lucas
correction, linear trace and an unshifted auxiliary representative.
The sign proof is essential: a positive coordinate chart cannot be
changed by an arbitrary integer shear. The unchanged large coefficient
tuple is supplied in full.

Two eleven-unrestricted-integer circuits retain exact degrees $298$ and
$300$ while reducing their costs to $146$ and $144$ binary gates.
These use a different complete coefficient tuple. No ten-positive result
is stated as an eleven-integer result or vice versa. Scalar support and
expanded coefficient height are not claimed when not materialized.
